\chapter{实验评估}
\label{ch:evaluation}
\chaptermark{实验评估}

本文测试构建／容量（RQ1）、执行（RQ2）、预算质量／复用（RQ3）和竞争区间（RQ4）。

\section{实验协议}
\label{sec:7-1}

\subsection{比较协议与可复现性}

\Cref{sec:3-2,sec:3-3,sec:3-5} 定义协议。搜索不计构建／校准；各消融使用自己的基线。

CPU JHQ 的设置／协议未匹配：十月 Vogue 使用 1,024 个列表、$[8,4]$ 位，而 E1 为 4,096 个列表。因此没有可匹配的 CPU／GPU 加速比。

自适应调用／仅设备计时见 \cref{sec:a-e1-protocol}。

指标与前沿估计器遵循 \cref{sec:3-2,sec:3-4}；$\glut$ 和 $\gscan$ 保持 \cref{sec:5-1} 中不同的含义。

\section{跨系统召回率与吞吐比较}
\label{sec:7-2}

\subsection{比较范围与前沿构造}

\Cref{fig:7-frontiers} 在 \cref{tab:3-data} 的前四个数据集上，将 JHQ-GPU 与 \cref{sec:3-4} 的图、压缩 IVF 和暴力搜索基线比较 \cite{cagra,pq,ivfrabitq}。E1 使用 batch 1,000 和共同 GT；T5 覆盖小 batch（\cref{sec:7-smallbatch}），内部消融见 \cref{sec:7-6}。

E1 将 105 个 arXiv／OpenAI 单元分派到 $\gscan=0$，27 个 Vogue 单元因变异系数（coefficient of variation，CV）为 0.338--0.622 分派到 $\gscan=4$。只有 Vogue 使用分组保留、共享阈值和归并（\cref{sec:5-5}）。

E1 有 1,590 个有效点：JHQ 每数据集 27--39 个，各基线家族 54--264 个。未出现失败、无效 ID 或重复行标记。参数扫描共享索引（\cref{sec:3-4}）。

\Cref{tab:3-grids,sec:b-grids} 定义有限覆盖，超出目标未测量。

存储不同：九位 RaBitQ 的名义编码字节为 $9d/8$；4/8 位码更少。\Cref{tab:4-bytes,sec:7-7} 在性能包络之外给出表示成本和实测分配。

\begin{figure}[!htbp]
\centering
\includegraphics[width=\linewidth]{ch7/frontiers.pdf}
\caption[E1 召回率—吞吐前沿]{四个 E1 数据集上的实测主机到主机召回率—QPS 前沿。显示范围从 0.90 开始，0.97 处画垂线。包括 CAGRA int8 在内的所有基线家族均参与前沿构造。线连接实测点；吞吐轴为对数。}
\label{fig:7-frontiers}
\end{figure}

\subsection{相同召回率下的吞吐}

\Cref{tab:7-matched} 使用 \cref{sec:3-4} 定义的插值和满足阈值的实测估计器报告召回率 0.97。

\par\medskip\noindent\begin{minipage}{\linewidth}
\captionsetup{hypcap=false}
\centering
\captionof{table}[召回率 0.97 时的跨系统比较]{召回率 0.97 的比较。每个比值为“插值／实测”；大于一有利于 JHQ。JHQ kQPS 为插值估计。arXiv 与 RaBitQ 的小幅差异随估计器改变方向。}
\label{tab:7-matched}
\small\setlength{\tabcolsep}{4pt}
% Generated by scripts/build_ch7.py; do not edit numbers manually.
\begin{tabular}{@{}lrrr@{}}
\toprule
数据集 & JHQ kQPS & 相对 CAGRA fp32 & 相对 IVF-RaBitQ \\
\midrule
Vogue & 230.8 & 1.416 / 1.537 & 1.059 / 1.046 \\
arXiv & 258.6 & 1.317 / 1.422 & 0.996 / 1.015 \\
OpenAI-1536 & 126.7 & 1.747 / 2.294 & 1.130 / 1.111 \\
OpenAI-3072 & 55.9 & 1.747 / 2.088 & 1.157 / 1.346 \\
\bottomrule
\end{tabular}

\end{minipage}\par\medskip

在索引、probes、预算及最终输出固定、batch 1,000 下，去掉分组使 Vogue 在召回率 0.97 时的 JHQ／RaBitQ 比从 1.059 降至 0.933。这分离的是集成分组路径，不只是共享阈值。

CAGRA int8 在 0.90--0.92 领先 $1.20$--$1.98\times$，但不覆盖 0.97。RaBitQ 赢得部分 0.98+ 目标；CAGRA fp32 在 arXiv 接近 0.99 时领先（\cref{fig:7-frontiers}）。

\newpage
\subsection{OpenAI 数据集上的精度敏感性}
\label{sec:7-precision}

补充源码检查确定，虽然输入／存储为 FP32，且源码注释有误导，E1 实现对 OpenAI 旋转和质心 GEMM 使用 TF32 乘法与 FP32 累加（\cref{sec:4-2}）。主 E1 前沿仍是该实现的历史测量。campaign D 中，独立 FP32 变体修改 cuBLAS 运算模式开关，在 batch 1,000 下重放 37/39 个 OpenAI 配置，每个五轮。不同 probe 设置间交替变体执行顺序。两者复用同一 TF32 训练模型；干预改变 add／查询算术，包括数据库编码与路由，并非比较独立 FP32 训练的模型。\Cref{tab:7-precision} 汇总成对变化。

\par\medskip\noindent\begin{minipage}{\linewidth}
\centering\captionsetup{hypcap=false}
\captionof{table}[匹配配置下的 FP32 敏感性]{FP32 敏感性。QPS 比为相同配置下 FP32／TF32 的几何平均；召回差采用同一方向。末两列比较召回率 0.97 下的 FP32 前沿与存档 E1 IVF-RaBitQ、CAGRA fp32 前沿。campaign D 未重测基线。}
\label{tab:7-precision}
\scriptsize\setlength{\tabcolsep}{3pt}
\begin{tabular}{@{}lrrrrr@{}}
\toprule
数据集 & 配置对 & QPS 比 & $\Delta$召回 & 相对 RaBitQ & 相对 CAGRA \\
\midrule
OpenAI-1536 & 37 & 0.953 & +0.0003--+0.0005 & 1.086 & 1.679 \\
OpenAI-3072 & 39 & 0.951 & +0.0000--+0.0003 & 1.079 & 1.629 \\
\bottomrule
\end{tabular}

\end{minipage}\par\medskip

相同配置下，FP32 吞吐降低约 5\%，召回增加 0--0.0005。召回率 0.97 时，FP32 比同会话 TF32 慢约 3\%，但仍为存档 RaBitQ 的 $1.086/1.079\times$、存档 CAGRA fp32 的 $1.679/1.629\times$。新会话 v69 对存档 RaBitQ 为 1.118/1.114，历史为 1.130/1.157。D 对基线的比较跨会话；仅 v69／v72 同会话。\Cref{fig:7-precision} 展示两条 JHQ 前沿与存档对照。T1 单独重测同会话锚点，网格敏感结果见 \cref{sec:b-precision-anchors}。

\par\medskip\noindent\begin{minipage}{\linewidth}
\captionsetup{hypcap=false}
\centering
\includegraphics[width=\linewidth]{supplement/precision.pdf}
\captionof{figure}[TF32 与 FP32 精度前沿]{campaign D 的精度前沿。实线为同会话 JHQ 变体；虚线／点线基线为历史 E1 测量。均使用本章前沿与插值规则。}
\label{fig:7-precision}
\end{minipage}\par\medskip

\subsection{同会话基线比较}
\label{sec:b-precision-anchors}
T1 使用 E1 内存池，每个 OpenAI 数据集十二个成对配置。基线运行前后各有一次五轮 JHQ 重复；中位数以几何方式组合。相同配置 FP32／TF32 均值为 0.956/0.964。两者均复用 TF32 训练缓存，改变 add／查询算术。
\par\medskip\noindent\begin{minipage}{\linewidth}
\captionsetup{hypcap=false}
\centering
\captionof{table}[同会话精度锚点]{T1，召回率 0.97、batch 1,000。基线比值为插值／严格实测阈值；FP32／TF32 为插值。所有方法处于同一会话，不同于 D 的存档基线。}
\label{tab:b-precision-anchors}
\footnotesize\setlength{\tabcolsep}{3pt}
\begin{tabular}{@{}lrrr@{}}
\toprule
数据集 & FP32/TF32 at 0.97 & FP32 / RaBitQ & FP32 / CAGRA \\
\midrule
OpenAI-1536 & 0.980 & 1.126 / 1.053 & 1.752 / 1.633 \\
OpenAI-3072 & 0.973 & 1.328 / 1.316 & 1.825 / 2.019 \\
\bottomrule
\end{tabular}

\end{minipage}\par\medskip
OpenAI-3072 的 RaBitQ 括区间 0.9669--0.9796 召回体现网格敏感性。这些锚点不替代 E1，也不提供 FP32 重训练结果。


\subsection{解释与比较限制}

两个 OpenAI 数据集均使用 Qdrant DBpedia 嵌入 \cite{data_openai1536,data_openai3072}。截断／重新归一化后，1,000 个查询行与 10,299 个抽样基础行的余弦为 $\geq0.9980$，中位数 1.000000，GT 重合度 9.046/10。没有实体 ID，这些检查支持行序对齐。维度翻倍后，CAGRA 比值仍为 1.747，RaBitQ 为 1.130/1.157，因此仅凭维度无法解释。

收益伴随更小的 128-probe 预算锚点：OpenAI 为 4--8，Vogue 为 64，arXiv 为 100，后续 FULL 为 200。这是相关观察，不是独立效应，也不是匹配前沿的参数设置。

\Cref{sec:7-8} 检查小幅领先的敏感性，以及 \cref{sec:3-3} 中接口、显存和编译差异的影响。

\subsection{小 Batch 改变竞争区间}
\label{sec:7-smallbatch}

T5 在四个 E1 数据集、batch 8/32 上比较调度 JHQ、CAGRA fp32／int8 和 IVF-RaBitQ。五轮同步主机到主机测量分别使用 1,000/992 查询。1,808 个测点包含 JHQ 单 block 对照，不含 IVF-PQ／暴力搜索。\Cref{fig:7-smallbatch,tab:7-smallbatch} 比较所得区间。

\suppressfloats[t]
\par\medskip\noindent\begin{minipage}{\linewidth}
\captionsetup{hypcap=false}
\centering
\includegraphics[width=\linewidth]{t4-t5/smallbatch.pdf}
\captionof{figure}[小 Batch 竞争区间]{T5 的 JHQ／最快对手 QPS 比，采用保留合格实测端点的插值。大于一有利于 JHQ。严格阈值见 \Cref{tab:7-smallbatch}。}
\label{fig:7-smallbatch}
\end{minipage}\par\medskip

\par\medskip\noindent\begin{minipage}{\linewidth}
\captionsetup{hypcap=false}
\centering
\captionof{table}[小 Batch 匹配召回比]{T5 QPS 比：保留合格端点的插值／严格实测阈值。各估计器独立选择最快对手；满足阈值的更高召回端点可入选，无需外推。}
\label{tab:7-smallbatch}
\footnotesize\setlength{\tabcolsep}{4pt}
\begin{tabular}{@{}lrrrr@{}}
\toprule
数据集 & $b$ & $R=0.90$ & $R=0.95$ & $R=0.97$ \\
\midrule
Vogue & 8 & 1.155 / 1.091 & 1.073 / 0.980 & 0.893 / 0.854 \\
Vogue & 32 & 1.069 / 0.997 & 0.956 / 0.895 & 0.865 / 0.839 \\
arXiv & 8 & 1.479 / 1.454 & 0.859 / 0.851 & 0.697 / 0.698 \\
arXiv & 32 & 1.166 / 1.148 & 0.667 / 0.661 & 0.622 / 0.582 \\
OpenAI-1536 & 8 & 1.316 / 1.131 & 0.851 / 0.825 & 0.641 / 0.620 \\
OpenAI-1536 & 32 & 1.071 / 0.915 & 0.770 / 0.752 & 0.937 / 0.914 \\
OpenAI-3072 & 8 & 1.032 / 0.953 & 0.798 / 0.798 & 0.721 / 0.702 \\
OpenAI-3072 & 32 & 0.866 / 0.734 & 0.608 / 0.771 & 1.008 / 1.351 \\
\bottomrule
\end{tabular}

\end{minipage}\par\medskip

batch 8／召回率 0.97 时，CAGRA fp32 在两种估计器下均赢得全部四个数据集，JHQ 比为 0.641--0.893。batch 32 时比为 0.622--1.008；实测阈值下 OpenAI-3072 有利于 JHQ。Vogue 的 batch-8／0.95 领先需要插值。网格分辨率会改变接近持平的结果。

纳入合格端点时，召回率 0.90 比值为 batch 8 的 1.032--1.479、32 的 0.866--1.166；严格实测下 JHQ 分别赢三／一个。省略端点会夸大收益。调度在 0.95 处相对单 block 提升 $1.7$--$3.3\times$，但不保证跨系统获胜。

空调用中位数为 0.012--0.072\,ms，仍计入总时间。CAGRA 使用 Python，JHQ／RaBitQ 使用 C++；预热不同（\cref{sec:a-t4-t5}）。


\section{Cartesian LUT、打包载入与 Lane 精排}
\label{sec:7-3}

RQ2 消融同时测量内核与应用效应。

\subsection{LUT 分解与吞吐}

LUT 消融测试 \cref{sec:5-2} 成本模型与 FP32 排序效应。

\Cref{fig:7-lut} 比较六个数据集上的 Cartesian 表分组，probe 数为 8/32/128/512，固定 $\alpha=100$，$\glut\in\{1,2,4,8\}$。五轮使用 batch 1,000，Stella 为 999。两组在 22/24 单元领先，例外接近持平：arXiv／128 的一组快 0.23\%，OpenAI-3072／8 的四组快 0.84\%。两组对一组成对获胜 23/24，全部数据集几何平均大于一，范围 1.021--1.318。共享索引单元并非独立试验。

\begin{figure}[!htbp]
\centering
\includegraphics[width=\linewidth]{ch7/lutgroups.pdf}
\caption[Cartesian 查找表分组]{固定 $\alpha=100$ 下的 Cartesian 表分组（v71）。组数表示 $\glut$，不是并行扫描组。曲线为各数据集相对一组实验臂 QPS 比的中位数；阴影表示范围。}
\label{fig:7-lut}
\end{figure}

两组／一组比值从 $0.998$ 到 $1.914$。四组和八组均存 16 项，但需要更多查表与加法。这说明表更小不一定提高吞吐；占用率和指令成本未被单独分离。任意四臂单元中报告的最大召回差为 $10^{-4}$。128 probes、$\alpha=100$ 下，T4 在 Vogue、arXiv 和 OpenAI-3072 上比较相同候选。整表与 $16+16$ 表有 22.6--26.1\% 的候选分数 bit 模式不同，最大相对误差 $4.53\times10^{-7}$。跨全部形式，20.5--29.4\% 不同，最大相对误差 $4.64\times10^{-7}$。Vogue 有 2--4/256 个最终 ID 差异，arXiv／OpenAI 没有。主码保留集合不变时，最终 ID 仍可变化，因为舍入后的主码分数进入精排（\cref{sec:a-t4-t5}）。

\subsection{打包访问与协作残差评分}

打包在 48 对匹配配置中获胜 $1.012$--$1.658\times$，覆盖六个数据集、四个 probe 数及 $\glut\in\{1,2\}$，召回差最多 $10^{-4}$。两个布局均为 subspace-major。打包／载入／解包一起改变。每臂仅一个汇总计时，比五轮 v71 的计时证据少。

\Cref{fig:7-packing-refinement,tab:7-refinement} 进一步区分残差精排中的宽载入和 lane 协作。L0 是原始标量路径；L1 只改变载入路径，不把一个候选分到多个 lanes；L16 为每候选分配 16 lanes。同一进程在五轮计时中轮换实验臂顺序，每轮每臂三次搜索。十二个普通单元中，L16 整体提升 $1.015$--$1.789\times$；独立的高预算 arXiv／$\alpha=200$ 对照达到 $2.383\times$。

\begin{figure}[!htbp]
\centering
\includegraphics[width=\linewidth]{ch7/packing_refinement.pdf}
\caption[打包访问与精排对照]{打包与精排对照。左：各表分组／probe 数下，六个数据集的打包载入 QPS 收益中位数及最小—最大阴影范围。右：128 probes、各数据集普通预算下，仅改变载入和16-lane 的实测端到端收益。两个面板使用不同实验家族。}
\label{fig:7-packing-refinement}
\end{figure}

\begin{table}[!htbp]
\centering
\caption[128 probes 下的精排收益]{128 probes 下的精排。比值为相对 L0 的成对端到端收益；末列除外，它是剖析得到的精排内核收益。内核与应用收益不能互换。}
\label{tab:7-refinement}
\small\setlength{\tabcolsep}{4pt}
% Generated by scripts/build_ch7.py; do not edit numbers manually.
\begin{tabular}{@{}lrrrr@{}}
\toprule
数据集 & $\alpha$ & $L1/L0$ QPS & $L16/L0$ QPS & 内核加速 \\
\midrule
arXiv & 100 & 1.368 & 1.789 & 5.97 \\
BGE & 32 & 1.026 & 1.062 & 1.97 \\
OpenAI-1536 & 8 & 1.055 & 1.145 & 3.34 \\
OpenAI-3072 & 4 & 1.063 & 1.171 & 8.81 \\
Stella & 16 & 1.030 & 1.053 & 2.66 \\
Vogue & 64 & 1.161 & 1.240 & 3.72 \\
\bottomrule
\end{tabular}

\end{table}

仅改变载入的 L1 与 L0 逐位相同，却在 arXiv 获得 $1.368\times$。在 $\alpha=100$ 下使用 lanes，精排从 4,277.5$\to$716.8\,$\mu$s，内核提升 $5.97\times$，而整体为 $1.789\times$。OpenAI-3072 的内核／整体提升分别为 $8.81\times$／$1.171\times$。其他阶段仍存在。

基线 wall time 占比为 $f$、精排加速比为 $s$ 时，替换模型给出
\[
 \frac{T'}T=(1-f)+\frac f s,\qquad
 \frac T{T'}=\frac1{1-f+f/s}.
\]
arXiv／OpenAI-3072 的 $f=0.529/0.168$ 预测 $1.787/1.175\times$，实测为 $1.789/1.171\times$。独立剖析提供解释性估计。

Lane FP32 重新结合运算的相对距离误差达到 $\approx1.6\times10^{-6}$，偶尔改变顺序。Sanitizer 只覆盖记录的子集，未声称普遍逐位相同。

\section{选择正确性与性能}
\label{sec:7-4}

\subsection{计算分数选择的严格 CPU 参考}
\label{sec:7-strict-selection}

全部 55,296 个分组查询实例的主码顺序与 CPU 参考精确一致。单 block 始终满足距离顺序，56 个集合差异均限于截断并列；全部 69,120 个实例对已保留精排分数的最终选择均精确（\cref{tab:7-strict-selection}）。

严格 CPU 检查对 GPU 生成的分数按 $(\primarydist,\mathrm{position})$ 排序。导出器遵循扫描布局和累加顺序，有效返回候选的导出分数逐位匹配。源码比较发现相对表扫描版本只增加导出。270 个配置覆盖 Vogue、arXiv 和 OpenAI-3072，probes 为 1、8、128，$\alpha\in\{4,100,200\}$、$\gscan\in\{0,1,2,4,8\}$ 以及两种阈值设置。

\suppressfloats[t]
\begin{table}[!htbp]
\centering
\caption[严格计算分数选择检查]{严格选择检查 CPU 参考检查。集合与顺序相等使用完整主码键。有效计数检查 $\min(c,n_c)$ 个保留候选；最终相等按 $(\refineddist,\mathrm{ID})$ 对保留的计算精排分数排序，包含其 FP32 bits。计数为查询实例，同一查询跨配置复用。}
\label{tab:7-strict-selection}
\footnotesize\setlength{\tabcolsep}{3pt}
\begin{tabular}{@{}lrrrrr@{}}
\toprule
路径 & 查询数 & 集合／顺序精确 & 仅并列 & 有效计数 & 最终精确 \\
\midrule
$\gscan\geq1$ & 55,296 & 55,296 / 55,296 & 0 & 55,296 & 55,296 \\
$\gscan=0$ & 13,824 & 13,768 / 8,758 & 56 & 13,824 & 13,824 \\
\bottomrule
\end{tabular}

\end{table}

全部 19,120 个候选不足案例的有效候选数正确。这些结果检验了 \cref{lem:5-buffer,prop:5-local-coverage,prop:5-shared-bound} 在记录案例中的表现。最终选择相等针对已保留的计算分数，而非原始空间精确近邻；sanitizer 覆盖限于其记录子集（\cref{sec:a-t4-t5}）。


\subsection{共享阈值保持选择目标}

48 组输出全部一致，QPS 比为 $0.994$--$1.562\times$。在 arXiv $\alpha=200$、batch 1,000、八组下，共享使归并从 2,629$\to$471\,$\mu$s，扫描从 7,968$\to$5,513\,$\mu$s。部分单元略有退步。

试验第一组包含 46--81\% 的优胜者。相同配额需要回退，不同配额仍留有瓶颈。阈值避免固定配额。

\subsection{探测列表选择与输出一致性}

Radix 保持评分算术／OpenAI TF32（\cref{sec:7-precision}）。全部 1,792 个合成用例通过。对 arXiv／Vogue／OpenAI-3072 的全部 1,000 个查询，probes 以及主码分数／位置均匹配。

\par\medskip\noindent\begin{minipage}{\linewidth}
\captionsetup{hypcap=false}
\centering
\captionof{table}[Probe 选择吞吐比]{精确计算分数 probe 改动的 v69／v65 端到端吞吐比。每个工作负载—batch 比较中，每个二进制有十轮计时，交替运行收集。$>1$ 的比值有利于 v69。}
\label{tab:7-probe}
\small
% Generated by scripts/build_ch7.py; do not edit numbers manually.
\begin{tabular}{@{}lrrrr@{}}
\toprule
数据集 & Batch 8 & 32 & 256 & 1,000 \\
\midrule
arXiv & 1.074 & 1.063 & 1.035 & 1.034 \\
Vogue & 1.109 & 1.103 & 1.052 & 1.041 \\
OpenAI-3072 & 1.167 & 1.117 & 1.055 & 1.050 \\
\bottomrule
\end{tabular}

\end{minipage}\par\medskip

\Cref{tab:7-probe} 显示端到端提升 $1.034$--$1.167\times$。batch 8 下，剖析的 probe 阶段在 arXiv 从 56.5 降至 22.6\,$\mu$s，Vogue 从 52.7 降至 14.7\,$\mu$s，OpenAI-3072 从 77.8 降至 16.7\,$\mu$s。未改阶段限制全 batch 收益，各比值使用自己的基线。

Vogue 在跨二进制的 1,000 查询转储中，probe／主码逐字节一致（一组／batch-8），最终召回却为 0.9642/0.9644。基线自身在一致性／计时运行间也由 0.9642 变为 0.9644。本次检查中，变化发生在主码选择之后，原因未解决。同进程内调度／单 block 输出匹配（\cref{app:reproducibility}）。

可选 FP16 probe 集合在 arXiv／Vogue／OpenAI-3072 的 1,000 查询中分别有 936/979/998 个与 v69 一致。高度重合不等于精确。这些结果不纳入主 E1 前沿，也不纳入精确 probe 主张。

\section{留出预算质量、吞吐与回本}
\label{sec:7-5}

RQ3 将 \cref{ch:refinement} 的策略与固定 $\alpha=100$ 及更强的通用对照 $\alpha=64$ 比较。

\subsection{无标签选择与有标签诊断 Oracle}

六个数据集在 128、512 probes 下测试 \cref{sec:6-4} 的固定、抽样、FULL、枚举和 oracle 对照。不相交查询池及每样本量 64 次抽样遵循 \cref{sec:3-2}。

所有规则变体均使用 \cref{eq:6-agreement}，取 $k=10$、$\epsilon=1$。FULL 使用全部校准查询；诊断 oracle 使用留出标签（\cref{tab:6-controls}）。

\subsection{全池校准作为质量优先工作点}

每个单元对全部 500 个校准查询执行一次全池校准。所选实验臂在 12/12 单元满足留出 $10^{-3}$ 损失目标，最大观测损失为 0.0006。其速率为 GT oracle 的 0.908--1.000，校准成本 9.5--60.3\,ms。相对固定 100 的速率几何平均比为 $1.111\times$；只看固定 100 也达标的十个单元，则为 $1.193\times$。\Cref{tab:7-budget-full} 给出完整比较。相对固定 64，全池选择在六个单元更快、两个持平、四个更慢；主要优势是观测中覆盖了全部十二个质量目标。

这些单元支持质量优先时用 FULL，校准成本优先时用 $S=128$（\cref{sec:6-2}）。

\begin{table}[!htbp]
\centering
\caption[全池预算校准]{全池校准（$S=500$）。损失相对最佳实测留出召回率；速率为所选固定臂的 kQPS。“Cal.”是每单元一次全池校准调用，不是重复查询流计时。留出查询为 500 个，Stella 为 499 个。}
\label{tab:7-budget-full}
\footnotesize\setlength{\tabcolsep}{3pt}
% Generated by scripts/build_ch7.py; do not edit numbers manually.
\begin{tabular}{@{}lrrrrrrr@{}}
\toprule
数据集 & Probe & $\alpha_{\rm full}$ & 损失 & $Q_{\rm full}$ & 相对 100 & 相对 oracle & 校准 ms \\
\midrule
Vogue & 128 & 100 & 0.0000 & 202.9 & 1.000 & 0.908 & 10.1 \\
arXiv & 128 & 200 & 0.0000 & 151.5 & 0.728 & 1.000 & 10.0 \\
OpenAI-1536 & 128 & 8 & 0.0000 & 267.7 & 1.599 & 1.000 & 9.5 \\
OpenAI-3072 & 128 & 4 & 0.0000 & 145.0 & 1.691 & 1.000 & 17.1 \\
BGE & 128 & 64 & 0.0006 & 124.9 & 1.075 & 0.932 & 17.1 \\
Stella & 128 & 16 & 0.0006 & 110.4 & 1.153 & 1.000 & 20.6 \\
Vogue & 512 & 100 & 0.0002 & 100.1 & 1.000 & 0.969 & 20.9 \\
arXiv & 512 & 200 & 0.0000 & 69.0 & 0.831 & 1.000 & 25.1 \\
OpenAI-1536 & 512 & 8 & 0.0000 & 114.9 & 1.326 & 1.000 & 21.8 \\
OpenAI-3072 & 512 & 4 & 0.0002 & 52.1 & 1.219 & 0.977 & 40.5 \\
BGE & 512 & 64 & 0.0006 & 50.8 & 1.033 & 1.000 & 42.3 \\
Stella & 512 & 16 & 0.0006 & 36.0 & 1.041 & 1.000 & 60.3 \\
\bottomrule
\end{tabular}

\end{table}

\subsection{低成本抽样及其质量风险}

\Cref{tab:7-budget-quality} 报告 $S=128$ 结果。十二个单元中，十个所选预算中位数匹配 oracle；arXiv／128 probes 和 OpenAI-3072／512 不同。预算跨工作负载为 4--200。预算中位数与平均召回不一定对应同一实验臂。平均召回由打印的最佳召回减去平均损失重建，约保留四位小数。

\begin{table}[!htbp]
\centering
\caption[抽样预算的留出质量]{$S=128$ 下的留出质量。规则 $\alpha$ 为 64 次抽样中位数；oracle $\alpha$ 使用留出标签。$\overline R_{\rm rule}$ 为重建的平均召回。“Losses”计数相对最佳已测留出召回损失 $>10^{-3}$ 的抽样次数；$L_{95}$ 为日志中的较低次序 95 分位损失估计。$R_{100}$ 为固定预算召回。}
\label{tab:7-budget-quality}
\small\setlength{\tabcolsep}{3pt}
% Generated by scripts/build_ch7.py; do not edit numbers manually.
\begin{tabular}{@{}lrrrrrr@{}}
\toprule
数据集 & Probe & $\alpha$: 规则/oracle & $R_{100}$ & $\overline{R}_{\rm rule}$ & 越阈数 & $L_{95}$ \\
\midrule
Vogue & 128 & 64 / 64 & 0.9680 & 0.9677 & 5/64 & 0.0022 \\
arXiv & 128 & 100 / 200 & 0.9748 & 0.9751 & 40/64 & 0.0024 \\
OpenAI-1536 & 128 & 8 / 8 & 0.9280 & 0.9279 & 4/64 & 0.0000 \\
OpenAI-3072 & 128 & 4 / 4 & 0.9252 & 0.9252 & 0/64 & 0.0000 \\
BGE & 128 & 32 / 32 & 0.9198 & 0.9190 & 1/64 & 0.0008 \\
Stella & 128 & 16 / 16 & 0.9892 & 0.9884 & 6/64 & 0.0030 \\
Vogue & 512 & 64 / 64 & 0.9916 & 0.9915 & 2/64 & 0.0002 \\
arXiv & 512 & 200 / 200 & 0.9878 & 0.9885 & 32/64 & 0.0016 \\
OpenAI-1536 & 512 & 8 / 8 & 0.9740 & 0.9738 & 4/64 & 0.0000 \\
OpenAI-3072 & 512 & 4 / 8 & 0.9704 & 0.9702 & 0/64 & 0.0002 \\
BGE & 512 & 64 / 64 & 0.9670 & 0.9661 & 24/64 & 0.0014 \\
Stella & 512 & 16 / 16 & 0.9954 & 0.9945 & 7/64 & 0.0030 \\
\bottomrule
\end{tabular}

\end{table}

arXiv 在 128/512 probes 下分别有 40/64、32/64 次损失越阈，尽管平均损失只有 0.0009。固定 100 损失 0.0012/0.0016：比默认更好未必达标。BGE／512 有 24 次越阈且 $L_{95}=0.0014$；Stella 达到 0.0030；OpenAI-3072 没有越阈（\cref{tab:7-budget-quality}）。四次损失可与报告 $L_{95}=0$ 同时出现（\cref{sec:a-budget-contract}）。

\begin{figure}[!htbp]
\centering
\includegraphics[width=0.82\linewidth]{ch7/budget_risk.pdf}
\caption[不同样本量下的留出风险]{不同样本量下的留出风险。纵轴为相对最佳已测预算，留出召回损失超过 $10^{-3}$ 的校准抽样百分比。因为 \cref{eq:6-agreement} 保持阈值为一，增加 $S$ 也会收紧允许的槽位不一致比例。}
\label{fig:7-budget-risk}
\end{figure}

\subsection{吞吐估计与校准经济性}

按 \cref{sec:a-budget-contract} 的平均所选臂估计，$S=128$ 在十个单元超过固定 100，幅度 $1.028$--$1.691\times$。arXiv 为 0.908/0.917，但平均召回更好，没有吞吐回本。

相对固定 64，抽样赢得 9/12 单元：几何平均 $1.054\times$，范围 $0.821$--$1.367\times$，OpenAI 为 $1.052$--$1.367\times$。两个固定对照均未满足 arXiv 质量，无法作统一质量匹配比较。

\begin{table}[!htbp]
\centering
\caption[抽样校准经济性]{$S=128$ 下的校准经济性。速率单位为 kQPS；规则速率是平均所选臂 QPS。比值分别比较固定 100 和 64。“Batches”用 500 查询、Stella 为 499，给出相对固定 100 的回本次数；“Never”表示无吞吐回本。质量列单独标明固定 100 是否满足留出损失目标。}
\label{tab:7-budget-cost}
\footnotesize\setlength{\tabcolsep}{3pt}
% Generated by scripts/build_ch7.py; do not edit numbers manually.
\begin{tabular}{@{}lrrrrrrrr@{}}
\toprule
数据集 & Probe & $Q_{100}$ & $\overline Q_{\rm rule}$ & 相对 100 & 相对 64 & 校准 ms & 批次数 & \shortstack{固定 100\\质量} \\
\midrule
Vogue & 128 & 202.9 & 224.4 & 1.106 & 1.004 & 4.0 & 17.0 & 达标 \\
arXiv & 128 & 207.9 & 188.8 & 0.908 & 0.821 & 4.2 & 不回本 & 未达标 \\
OpenAI-1536 & 128 & 167.5 & 268.5 & 1.603 & 1.367 & 3.9 & 3.5 & 达标 \\
OpenAI-3072 & 128 & 85.7 & 145.0 & 1.691 & 1.230 & 5.7 & 2.4 & 达标 \\
BGE & 128 & 116.2 & 131.2 & 1.130 & 1.051 & 6.5 & 13.2 & 达标 \\
Stella & 128 & 95.7 & 110.3 & 1.153 & 1.087 & 7.5 & 10.8 & 达标 \\
Vogue & 512 & 100.1 & 102.9 & 1.028 & 0.996 & 7.5 & 55.4 & 达标 \\
arXiv & 512 & 83.1 & 76.2 & 0.917 & 0.870 & 9.7 & 不回本 & 未达标 \\
OpenAI-1536 & 512 & 86.6 & 115.0 & 1.327 & 1.268 & 8.8 & 6.2 & 达标 \\
OpenAI-3072 & 512 & 42.8 & 52.1 & 1.219 & 1.052 & 13.1 & 6.2 & 达标 \\
BGE & 512 & 49.2 & 51.1 & 1.039 & 1.006 & 15.8 & 41.2 & 达标 \\
Stella & 512 & 34.6 & 36.0 & 1.041 & 1.028 & 21.3 & 37.1 & 达标 \\
\bottomrule
\end{tabular}

\end{table}

校准成本 3.9--21.3\,ms。相对固定 100/64，最大正收益预计回本时间为 0.54/2.72\,s，FULL 为 1.52/2.25\,s（\cref{sec:6-5,eq:6-payback,tab:7-budget-cost}）。arXiv 固定对照质量不达标；下面的直接复用同时测成本和召回。

\subsection{完整决策轨迹与有限池风险}
\label{sec:7-budget-trace}

全部 2,304 条抽样曲线和十二条 FULL 曲线均单调，选择与枚举一致。记录的子集和复现 18,432 个网格计数（\cref{eq:6-finite-pool}）。这验证已记录汇总；任意子集／位置不变性及一般的仅距离嵌套仍有条件（\cref{sec:6-3}）。

$S=128$ 下最大损失为 arXiv／512 的 0.0050、BGE／512 的 0.0046，$S=32$ 下为 0.0294。FULL 保持在 0.0006 内。

arXiv／128 在 $S=32/64/128$ 下预测／观测越阈分别为 0.965/0.984、0.875/0.875、0.620/0.625（\cref{fig:7-finite-pool}）。36 个差距中最大为 BGE／512 的 0.144，即 0.519/0.375。GT 损失预测需要稳定计数／留出臂质量；校准本身不使用标签。

\begin{figure}[!htbp]
\centering
\includegraphics[width=\linewidth]{supplement/finite-pool.pdf}
\caption[有限池风险预测与观测]{有限池预测与观测到的 arXiv 损失越阈。误差条为固定池条件下 64 个独立抽样子集的一倍估计标准误，不表示未来工作负载的不确定性。}
\label{fig:7-finite-pool}
\end{figure}

十次交替重复支持低预算计时平台，而非一般运行时间单调性。

\subsection{直接测量校准加复用}
\label{sec:7-stream-cost}

四个单元在五轮轮换顺序中，对 500 个重复留出 batches 比较固定 64/100、抽样和 FULL。计时包含查询收集、校准、预算应用和首次图，不含构建、预热及扫描校准（\cref{tab:7-stream-cost}）。

\par\medskip\noindent\begin{minipage}{\linewidth}
\centering\captionsetup{hypcap=false}
\captionof{table}[500 个复用 Batch 的直接成本]{500 个复用 batches 的直接总毫秒数，取五轮平均，抽样与 FULL 策略包含预算校准。每 batch 有 500 查询。质量与计时结果一同讨论。}
\label{tab:7-stream-cost}
\footnotesize\setlength{\tabcolsep}{4pt}
\begin{tabular}{@{}lrrrr@{}}
\toprule
数据集／probes & 固定 100 & 固定 64 & 规则 $S=128$ & FULL \\
\midrule
arXiv/128 & 1224 & 1102 & 1291 & 1680 \\
OpenAI-1536/128 & 1491 & 1264 & 852 & 860 \\
OpenAI-3072/128 & 2757 & 2160 & 1769 & 1783 \\
BGE/512 & 5185 & 5018 & 4921 & 5062 \\
\bottomrule
\end{tabular}

\end{minipage}\par\medskip

OpenAI-3072 两种策略均选 4，召回 0.9252，与固定 100 相同。OpenAI-1536 的 FULL／固定 100 为 0.9280，抽样为 0.9276。时间见 \Cref{tab:7-stream-cost}。

相对固定 100/64，直接节省分别从 OpenAI-1536 的 4/6 个 batches、OpenAI-3072 的 4/9 个 batches 开始（\cref{fig:7-stream-cost}）；分数批次拟合结果有所不同。BGE 抽样／固定 100／FULL 召回为 0.9658/0.9670/0.9664。arXiv FULL 为质量选择更慢的 200，固定对照不达标。

\par\medskip\noindent\begin{minipage}{\linewidth}
\captionsetup{hypcap=false}
\centering
\includegraphics[width=\linewidth]{supplement/stream-cost.pdf}
\captionof{figure}[校准后复用的直接累计时间]{直接测得的平均累计时间，包含预算校准和首 batch 准备。每条曲线重放一个固定留出查询池。不同策略可能有不同召回率；仅靠时间不能形成质量匹配排名。数据集／probe 设置与 \cref{tab:7-stream-cost} 相同。}
\label{fig:7-stream-cost}
\end{minipage}\par\medskip


\section{批量大小、调度与负载不均}
\label{sec:7-6}

\subsection{并行度与重复选择工作}

batch-8 的八组提升 $2.06$--$2.76\times$；arXiv 扫描从 849$\to$216\,$\mu$s，另加 16.7\,$\mu$s 归并。batch 1,000、$\alpha=200$ 时，独立分组反而将扫描从 3,530$\to$7,652\,$\mu$s，另加 2,624\,$\mu$s 归并，仅达到单 block QPS 的 $0.46\times$。共享阈值减少工作，但未消除区间变化。

\begin{figure}[!htbp]
\centering
\includegraphics[width=\linewidth]{ch7/batch_execution.pdf}
\caption[依赖 Batch 的执行区间]{执行区间。左：固定共享阈值组数及调度执行，相对 lane 精排单 block 基线的速率；小 batch 有 13 单元，batch 1,000 有 11。右：arXiv 预算 100、200 在 batch 8、1,000 下的成对扫描／归并剖析。斜纹表示归并。右侧来自独立阈值扫描实验，解释一种成本，不是调度器的完整加速。}
\label{fig:7-batch}
\end{figure}

\subsection{在更多配置上评估固定规则}

调度遵循 \cref{sec:5-5} 的填充／偏斜规则与伪查询校准。四个配置先确定常数，再测量其他配置。

\begin{table}[!htbp]
\centering
\caption[集成执行与调度收益]{相对 v61 L16 的集成执行收益。条目为五轮逐单元成对加速比中位数的几何平均。“Held out”排除确定规则的四个配置，表示留出配置，而非预算实验的留出查询划分。}
\label{tab:7-scheduling}
\small\setlength{\tabcolsep}{4pt}
% Generated by scripts/build_ch7.py; do not edit numbers manually.
\begin{tabular}{@{}rrrrrrr@{}}
\toprule
Batch & 单元数 & $g_s=2$ & $g_s=4$ & $g_s=8$ & 调度 & 留出配置 \\
\midrule
8 & 13 & 1.318 & 2.014 & 2.728 & 2.727 & 2.930 \\
32 & 13 & 1.323 & 1.970 & 2.337 & 2.356 & 2.560 \\
256 & 13 & 0.997 & 1.090 & 1.099 & 1.057 & 1.041 \\
1000 & 11 & 0.959 & 0.947 & 0.886 & 1.015 & 1.012 \\
\bottomrule
\end{tabular}

\end{table}

平均值使用十三个小 batch 配置、十一个全 batch 配置；Stella 有 999 查询。留出子集分别为九／七个（\cref{tab:7-scheduling}）。batch 1,000 下，八组为 $0.886\times$；最坏 regret 为 1.217（BGE／256）。

Vogue 的 CV 0.338--0.622 在 batch 1,000 时选择四组。比值范围 0.895--1.175；$(n_{\rm probe},\alpha)=(16,16),(16,64),(32,64)$ 对应 0.955、0.895、0.955。CV 测量离散程度，而非总工作量：它使 0.97 处领先，却在低 probe 负载上过度分组（\cref{sec:7-2}）。

250 个变体／配置对、234,320 次查询比较中，没有非并列主码差异或最终 ID 集合差异。并列输出顺序可能不同。这些是顺序同步吞吐测试，不是开放到达流尾延迟测试。

\section{冷构建、常驻显存与已测规模}
\label{sec:7-7}

冷构建、常驻显存和完成的大数据集搜索用于评估 RQ1；E1 显存比较也支持 RQ4。

\subsection{冷构建与生命周期边界}

三次空缓存重复得到 Vogue 训练 298--303\,ms，总计 $\approx0.56$\,s，而较慢存档为 1.33\,s。Stella 为 10.54--10.65\,s，存档为 10.40\,s（\cref{tab:7-build-memory}）。

\begin{table}[!htbp]
\centering
\caption[冷构建与常驻显存]{冷 JHQ 构建与单独测量的常驻显存。Train／add 为三次补充正常构建的组件中位数；Total 为成对总和中位数。显存来自 v71 的两组、128 probes、$\alpha=100$ 运行，包含上下文和工作区。“Residual”是残差编码模型大小占实测使用总量的比例，不是每查询读取比例。}
\label{tab:7-build-memory}
\small\setlength{\tabcolsep}{3pt}
% Generated by scripts/build_ch7.py; do not edit numbers manually.
\begin{tabular}{@{}lrrrrrr@{}}
\toprule
数据集 & 训练 s & Add s & 合计 s & 使用 GiB & 原始 GiB & 残差占比 \\
\midrule
Vogue & 0.30 & 0.26 & 0.56 & 1.47 & 2.67 & 45.4\% \\
arXiv & 0.48 & 0.51 & 0.99 & 2.58 & 6.45 & 62.5\% \\
OpenAI-1536 & 0.61 & 0.38 & 1.00 & 2.39 & 5.72 & 59.7\% \\
OpenAI-3072 & 1.16 & 0.67 & 1.82 & 4.14 & 11.43 & 69.0\% \\
BGE & 3.56 & 3.65 & 7.20 & 11.93 & 38.50 & 80.7\% \\
Stella & 3.55 & 6.99 & 10.55 & 20.44 & 67.81 & 82.9\% \\
\bottomrule
\end{tabular}

\end{table}

\Cref{sec:3-3,sec:a-e1-protocol} 定义启动／构建边界；\cref{fig:7-build} 比较重复构建。

\par\medskip\noindent\begin{minipage}{\linewidth}
\captionsetup{hypcap=false}
\centering
\includegraphics[width=\linewidth]{ch7/construction.pdf}
\captionof{figure}[记录的构建成本]{记录的构建成本，纳入 CAGRA int8。JHQ 为输入页已热时三次补充正常冷构建的中位数及范围；其他柱为成功测量索引配置的中位数，误差线为最小—最大值。Int8 包含数据集量化；RaBitQ 不含所需序列化往返。缺失大数据集柱表示未尝试，不是实测构建失败。纵轴为对数。}
\label{fig:7-build}
\end{minipage}\par\medskip

Stella 的隔离 CAGRA int8 为 74.55/82.45\,s，含量化；三个 IVF-PQ 配置的中位数为 86.59\,s，而 JHQ 为 10.54--10.65\,s。质量／常驻显存另行比较。

五次 arXiv 构建的召回范围为 0.0012。Stella 启动在热页下为 18.2--18.6\,s，首次读取为 129.9\,s。20\,ms 采样峰值为 BGE 26.08\,GiB、Stella 22.58\,GiB。BGE 的 12\,GiB 分配缓冲先于 23.05\,GiB gather 载荷存活，不能相加（\cref{eq:4-assignment-scratch}）。

\subsection{常驻分配与压缩的作用}

\Cref{fig:7-memory} 将 \cref{sec:4-4} 组件模型与实测常驻总量结合，保留未归属余量。前述构建测量给出采样峰值。

\par\medskip\noindent\begin{minipage}{\linewidth}
\captionsetup{hypcap=false}
\centering
\includegraphics[width=\linewidth]{ch7/memory.pdf}
\captionof{figure}[常驻显存分解]{两组表、128 probes、$\alpha=100$ 下的常驻显存分解。组件大小为模型值，标记为实测使用总量，斜纹为未归属余量。总量包含上下文和工作区。前述构建测量给出采样峰值；这些常驻标记不是峰值。}
\label{fig:7-memory}
\end{minipage}\par\medskip

残差占 BGE 11.93\,GiB 总量的 80.7\%，占 Stella 20.44\,GiB 的 82.9\%，总量含上下文／工作区。原向量需要 38.50/67.81\,GiB。更小 $\alpha$ 节省读取，不减少存储残差 bytes（\cref{sec:4-4}）。

E1 测量在 \cref{tab:7-e1-memory} 提供独立跨系统视角。从 JHQ 进程总量减去基线测得的 507.2\,MiB 上下文，约为 CAGRA fp32 索引分配的 0.29--0.39。JHQ 包含工作区，CAGRA 则为构建时索引增量。这些是参数扫描范围，不是显存匹配召回点或等峰值预算比较。同表 CAGRA int8 与 JHQ 分配同量级：相对 fp32 的压缩并非 JHQ 独有。\cref{tab:4-bytes} 组件模型解释这一比较。

\par\medskip\noindent\begin{minipage}{\linewidth}
\captionsetup{hypcap=false}
\centering
\captionof{table}[E1 分配范围]{E1 分配范围。JHQ 为进程总量减去估计的 507.2\,MiB 上下文，包含查询工作区；两个 CAGRA 列是构建时测得的索引分配。范围覆盖参数扫描。}
\label{tab:7-e1-memory}
\footnotesize\setlength{\tabcolsep}{4pt}
% Generated by scripts/build_ch7.py; do not edit numbers manually.
\begin{tabular}{@{}lrrr@{}}
\toprule
数据集 & \shortstack{JHQ 净值\\进程 GiB} & \shortstack{CAGRA fp32\\索引 GiB} & \shortstack{CAGRA int8\\索引 GiB} \\
\midrule
Vogue & 0.90--1.01 & 2.89--2.91 & 0.89--0.91 \\
arXiv & 2.01--2.19 & 6.84--6.98 & 2.01--2.15 \\
OpenAI-1536 & 1.79--2.33 & 5.95--5.97 & 1.67--1.68 \\
OpenAI-3072 & 3.54--3.82 & 11.67--11.68 & 3.10--3.11 \\
\bottomrule
\end{tabular}

\end{minipage}\par\medskip

\subsection{大数据集可行性与竞争区间}

BGE 包含 10,091,524 个向量，Stella 为 17,776,615；二者均为 1,024 维、32,768 个 IVF 列表。Stella 查询从语料中抽取，因此结果刻画这一基准，不是独立采集的外部查询分布。\Cref{fig:7-large} 纳入成功的 CAGRA int8 与 IVF-PQ 测量。预计存储超容量时，完整精度 CAGRA／RaBitQ／FP32 暴力搜索被预先排除。这是计划排除，不是实测 OOM 或召回失败。

\par\medskip\noindent\begin{minipage}{\linewidth}
\captionsetup{hypcap=false}
\centering
\includegraphics[width=\linewidth]{ch7/large_frontiers.pdf}
\captionof{figure}[大数据集搜索前沿]{大数据集主机到主机前沿，纳入成功的隔离 CAGRA int8 重跑。BGE／Stella 使用 1,000/999 查询。端点标记已测网格覆盖。}
\label{fig:7-large}
\end{minipage}\par\medskip

CAGRA int8 可容纳 Stella。在 BGE 召回率 0.88--0.95 上，它领先 $1.35$--$2.96\times$；在 int8 最高召回率 0.9579 处，JHQ 领先 $1.086\times$。BGE／Stella 最高召回分别为：int8 0.9579/0.9783，JHQ 0.9663/0.9952，PQ 0.9525/0.9896。

Stella 上，JHQ 在召回 0.978679 时为 185.4 kQPS，int8 在 0.978278 时为 66.3。$2.80\times$ 的测点比并未匹配召回；int8 不覆盖更高区间（\cref{fig:7-large}）。

\section{数值检查与召回分析}
\label{sec:7-8}

\subsection{质量分辨率与重复执行}

一个 GT 命中使 Recall@10 改变
\[
 \Delta r=\frac1{nk}=
 \begin{cases}10^{-4},&n=1000,\\2\times10^{-4},&n=500.\end{cases}
\]
Stella 的 $n=499$ 步长略大。这个步长不是不确定区间。

全部 369 个 E1 基线重复 ID 标记均保留，包括并列重排。401 口径增加跨协议差异（\cref{sec:a-e1-protocol}）；这些不是独立构建。

0.97 附近，吸附前的局部斜率对应每 $10^{-3}$ 召回变化 1.01--3.17\% QPS。五次构建的召回范围在 Vogue／arXiv／OpenAI-1536/3072 上为 0.0008/0.0012/0.0004/0.0002。这提示局部敏感性，不是重新训练前沿的界。缓存模型 FP32 检查保持 OpenAI 优势（\cref{sec:7-precision}）；arXiv／LUT 的接近持平仍无定论。

RaBitQ 在 Vogue／arXiv／OpenAI-1536/3072 上，五轮范围占中位 QPS 的中位数为 3.79\%、3.34\%、2.87\%、5.25\%。这些离散范围不是置信区间。为降低偶发 30--60\,ms 调用的影响，本文将 1,000 查询速率重算为 $10^6/\operatorname{median}(\text{call ms})$。召回率 0.97 时，所得 JHQ／RaBitQ 比为 1.063、0.995、1.118、1.100。三个有利方向和 arXiv 接近持平仍保持，因此结果并非仅由长调用造成。

\par\medskip\noindent\begin{minipage}{\linewidth}
\setlength{\parindent}{2em}
\subsection{重建评分与召回损失}
\label{sec:7-recall-attribution}

\cref{tab:5-query-stages} 的阶段约束区分计算分数选择与原始空间 kNN；lane 舍入、cutoff 并列和进程变化分别限定不同阶段。

存档重建试验测量距全 probe 重建评分的剩余差距。七个工作负载—预算臂各有 128 查询，对全部 probed 候选做 FP64 评分，相比所选预算的 GPU 输出，召回提高 0--0.002344。arXiv／128 probes 下，$\alpha=200$ 在全部 128 查询上匹配全 IVF 重建答案集。预算为 100 时，召回 0.9734375，全 IVF 为 0.9750；在路由／重建共同缺口 0.0250 之外，还有 0.0015625 的预算缺口（\cref{sec:6-2}）。这些汇总包含同一数据集—probe 工作负载的两个预算。FP64／GPU 算术差异使部分实验臂不能纯粹归因为预算。分离路由／表示需要逐阶段 GT 覆盖。
\end{minipage}\par
